\documentclass[12pt,a4paper]{article}

\usepackage[T1]{fontenc}
\usepackage[utf8]{inputenc}
\usepackage[brazil]{babel}
\usepackage{lmodern}
\usepackage{geometry}
\usepackage{hyperref}
\usepackage{enumitem}
\usepackage{array}

\geometry{margin=2.5cm}
\hypersetup{
  colorlinks=true,
  linkcolor=blue,
  urlcolor=blue
}

\title{Plano de Aula: Aula 7 -- Teste de Software}
\date{}

\begin{document}
\maketitle

\noindent
\textbf{Disciplina:} Engenharia de Software (Curso Técnico em Desenvolvimento de Sistemas)\\
\textbf{Duração:} 2 horas (120 minutos)\\
\textbf{Modalidade:} Presencial\\
\textbf{Materiais da aula:} \href{aula-07.html}{slides (HTML)} \quad \href{aula-07.pdf}{ementa (PDF)}\\
\textbf{Livro-base (referência):} \href{../99-Materiais/Engenharia%20de%20software%20projetos%20e%20processos%20by%20Wilson%20de%20P%C3%A1dua%20Paula%20Filho%20.epub.pdf}{Engenharia de software: projetos e processos (PDF)}

\section{Competências e Habilidades}
\textbf{Competências}
\begin{itemize}[leftmargin=*,itemsep=0.25em]
  \item Entender e implementar os principais tipos de teste de software.
  \item Compreender qualidade como parte do processo de desenvolvimento.
\end{itemize}

\textbf{Habilidades}
\begin{itemize}[leftmargin=*,itemsep=0.25em]
  \item Diferenciar níveis de teste e objetivos.
  \item Criar casos de teste a partir de requisitos e critérios de aceitação.
  \item Identificar cenários de erro, limite e exceções.
\end{itemize}

\section{Objetivos da Aula}
\begin{itemize}[leftmargin=*,itemsep=0.25em]
  \item Definir teste e diferenciar defeito, falha e causa.
  \item Entender níveis de teste: unitário, integração, sistema e aceitação.
  \item Escrever casos de teste completos com pré-condições, passos e resultado esperado.
  \item Discutir estratégia mínima de testes e a ideia da pirâmide de testes.
\end{itemize}

\section{Assuntos Abordados no Dia}
\begin{itemize}[leftmargin=*,itemsep=0.25em]
  \item Conceitos: teste, defeito, falha.
  \item Níveis de teste e quando usar cada um.
  \item Caso de teste: estrutura e boas práticas.
  \item Pirâmide de testes e equilíbrio entre velocidade e cobertura.
  \item Atividade: derivar testes a partir de requisitos do ``app organizador''.
\end{itemize}

\section{Metodologia}
Aula expositiva com exemplos curtos e ligação direta com requisitos (Aula 4). Exercício em grupos para escrever casos de teste e identificar cenários de erro/limite. Fechamento com discussão sobre estratégia de testes e trade-offs.

\section{Materiais e Recursos}
\begin{itemize}[leftmargin=*,itemsep=0.25em]
  \item Quadro e marcadores.
  \item Folhas A4 para escrita de casos de teste.
  \item Slides: \href{aula-07.html}{aula-07.html} e ementa: \href{aula-07.pdf}{aula-07.pdf}.
  \item Livro-base (PDF): \href{../99-Materiais/Engenharia%20de%20software%20projetos%20e%20processos%20by%20Wilson%20de%20P%C3%A1dua%20Paula%20Filho%20.epub.pdf}{Engenharia de software: projetos e processos}.
  \item Vídeos complementares (links nos slides).
\end{itemize}

\section{Cronograma e Descrição Detalhada (120 Minutos)}
\subsection*{Parte 1: Abertura e motivação (10 min)}
\begin{itemize}[leftmargin=*,itemsep=0.35em]
  \item Retomar critérios de aceitação (Aula 4) e conectá-los aos testes.
  \item Apresentar objetivos e entregas da atividade.
\end{itemize}

\subsection*{Parte 2: Conceitos e níveis de teste (35 min)}
\begin{itemize}[leftmargin=*,itemsep=0.35em]
  \item Definir teste e diferenciar defeito/falha.
  \item Níveis de teste com exemplos do ``app organizador''.
  \item Discussão: por que testar cedo reduz custo.
\end{itemize}

\subsection*{Parte 3: Como escrever casos de teste (25 min)}
\begin{itemize}[leftmargin=*,itemsep=0.35em]
  \item Estrutura: pré-condições, entradas, passos, resultado esperado.
  \item Exemplos: cenário de sucesso e cenário de erro.
  \item Checklist: reprodutível, claro, ligado a um requisito.
\end{itemize}

\subsection*{Parte 4: Atividade prática (35 min)}
\begin{itemize}[leftmargin=*,itemsep=0.35em]
  \item Em grupos, selecionar 2 requisitos do backlog.
  \item Para cada requisito:
  \begin{itemize}[leftmargin=*,itemsep=0.25em]
    \item Escrever 2 casos de teste (incluindo 1 caso de erro/limite).
    \item Identificar nível de teste (unitário/integração/sistema/aceitação).
  \end{itemize}
  \item Troca entre grupos: revisar um caso de teste de outro grupo e sugerir melhoria.
\end{itemize}

\subsection*{Parte 5: Estratégia e fechamento (15 min)}
\begin{itemize}[leftmargin=*,itemsep=0.35em]
  \item Explicar pirâmide de testes e equilíbrio: muitos testes caros deixam o ciclo lento.
  \item Ponte para Aula 8: planejamento e gerenciamento também reduzem risco e aumentam qualidade.
\end{itemize}

\section{Avaliação}
\begin{itemize}[leftmargin=*,itemsep=0.25em]
  \item Qualidade dos casos de teste (clareza, completude e verificabilidade).
  \item Cobertura de cenários de erro/limite e coerência com requisitos.
  \item Participação na revisão cruzada.
\end{itemize}

\section{Referências}
\begin{itemize}[leftmargin=*,itemsep=0.25em]
  \item \href{../99-Materiais/Engenharia%20de%20software%20projetos%20e%20processos%20by%20Wilson%20de%20P%C3%A1dua%20Paula%20Filho%20.epub.pdf}{PAULA FILHO, Wilson de Pádua. Engenharia de software: projetos e processos. 4. ed. 2019.}
\end{itemize}

\end{document}
